% ----------------------------------------------------------------------
% Capítulo 4: Diagnóstico e validação do produto
% ----------------------------------------------------------------------
\chapter{Diagnóstico e validação do produto}
\label{cap:diagnostico}

\instrucao{De oito a quinze páginas. Primeiro os resultados do diagnóstico, depois a descrição do produto e, por fim, a validação. Cada tabela ou figura é citada no texto antes de aparecer.}

\section{Diagnóstico das bibliotecas}
A Tabela~\ref{tab:praticas} resume as práticas declaradas pelas equipes (dados de exemplo).

\begin{table}[htb]
  \centering
  \caption{Práticas de formação de acervo declaradas ($n = 12$)}
  \label{tab:praticas}
  \begin{tabular}{lrr}
    \toprule
    \textbf{Prática} & \textbf{$f_i$} & \textbf{\%} \\
    \midrule
    Possui critérios escritos de seleção & 2 & 16,7 \\
    Registra as doações recebidas & 7 & 58,3 \\
    Fez descarte nos últimos três anos & 3 & 25,0 \\
    Consulta os leitores antes de comprar & 5 & 41,7 \\
    \bottomrule
  \end{tabular}
  \fonte{Dados da pesquisa (exemplo).}
\end{table}

\instrucao{Interprete a tabela em vez de repeti-la. Exemplo: o dado mais baixo aponta a seção do guia que precisa de mais detalhe.}
O dado que mais chama atenção é a ausência de critérios escritos em dez das doze bibliotecas. Esse resultado definiu a seção de seleção como a parte mais detalhada do guia.

\section{Descrição do produto}
\instrucao{Apresente o produto: público, formato, seções e modo de uso. O produto completo vai no apêndice; aqui entra a visão geral.}

O guia tem cinco seções, uma para cada elemento do Quadro~\ref{qua:elementos}, e quatro formulários. Cada seção traz um texto-modelo que a escola completa com seus dados e aprova no conselho escolar.

\section{Validação}
\instrucao{Relate como o produto foi testado, quem avaliou, quais critérios foram usados e o que mudou na versão final.}

As oficinas reuniram as equipes das doze bibliotecas. As sugestões mais frequentes foram simplificar o formulário de doação e incluir exemplos de justificativa de descarte, ambas incorporadas à versão final.
